\subsubsection{Dinic (Max Flow)}
Dinic 算法，用于高效求解网络最大流。
\begin{lstlisting}
struct node{ll u,v,w,nx;}e[1000010];
ll hd[100010],cnt=1;
void add(ll u,ll v,ll w){e[++cnt]={u,v,w,hd[u]};hd[u]=cnt;}
void ad(ll u,ll v,ll w){add(u,v,w),add(v,u,0);}

ll de[100010],cur[100010];
ll dis[100010];
const ll inf=LLONG_MAX/4;
queue<ll>q;
ll nodes,S,T;
void init_dinic(int count,int source,int sink){
    nodes=count;S=source;T=sink;cnt=1;fill(hd,hd+nodes+1,0);
}
bool bfs(){
    while(!q.empty())q.pop();
    f(i,1,nodes)cur[i]=hd[i],de[i]=inf;de[S]=0;q.push(S);
    while(!q.empty()){
        ll u=q.front();q.pop();
        for(int i=hd[u];i;i=e[i].nx){
            ll v=e[i].v;if(e[i].w==0)continue;
            if(de[v]==inf){
                de[v]=de[u]+1;q.push(v);
                if(v==T)return 1;
            }
        }
    }return 0;
}
ll dfs(ll u,ll fl){
    if(u==T)return fl;
    ll res=0,k=0;
    for(ll &i=cur[u];i&&fl;i=e[i].nx){if(e[i].w==0)continue;
        ll v=e[i].v;if(de[v]==de[u]+1){
            k=dfs(v,min(fl,e[i].w));
            e[i].w-=k,e[i^1].w+=k;
            fl-=k,res+=k;
            if(!fl)break; // preserve a partially usable current arc
        }
    }return res;
}
ll n,m;
ll dinic(){
    ll res=0;
    if(S==T)return 0;
    while(bfs()){ll pushed;while((pushed=dfs(S,inf)))res+=pushed;}
    return res;
}
int main(){
    n=d,m=d,S=d,T=d;init_dinic(n,S,T);
    f(i,1,m){
        ll u=d,v=d,w=d;
        ad(u,v,w);
    }cout<<dinic();
    return 0;
}
\end{lstlisting}


\paragraph{使用条件与多测} 每组读边前清 hd[0..nodes]、cnt=1；容量残量由 ad 重新赋值。重复 dinic() 返回新增流量而非历史总流量。容量非负，源汇应不同，流量和必须可用 ll 表示；nodes 包括超级源汇，边池容量按两倍有向边数算。DFS 可能深至点数；通用图复杂度 $O(V^2E)$，不能将单位容量图的更强界套到任意网络。
